% 总变差的分割读法：X={1,2,3,4}，eta({i})=2,-1,1,-2。三种分割的sum|eta(A_i)|分别为|0|=0、|1|+|-1|=2、|3|+|-3|=6；|eta|(X)=6。
\begin{tikzpicture}[
  lab/.style={font=\figurefont\fontsize{8}{10}\selectfont,text=FigureInk},
  note/.style={font=\FigureNoteFont\fontsize{7.5}{9}\selectfont,text=FigureMuted,align=center},
  pt/.style={circle,draw=FigureStroke,line width=.7pt,minimum size=6.2mm,inner sep=0pt,font=\figurefont\fontsize{8}{10}\selectfont},
  ppt/.style={pt,fill=FigureBlueFill},
  npt/.style={pt,fill=FigureRedFill},
  grp/.style={draw=FigureMuted,line width=.7pt,rounded corners=4pt}]
\foreach \r/\y in {1/0,2/-1.1,3/-2.55}{
  \begin{scope}[yshift=\y cm]
    \node[ppt] (a\r) at (0,0) {$+2$};
    \node[npt] (b\r) at (.9,0) {$-1$};
    \node[ppt] (c\r) at (1.8,0) {$+1$};
    \node[npt] (d\r) at (2.7,0) {$-2$};
  \end{scope}}
\node[lab] at (0,.6) {$1$}; \node[lab] at (.9,.6) {$2$}; \node[lab] at (1.8,.6) {$3$}; \node[lab] at (2.7,.6) {$4$};
% row 1: whole set
\draw[grp] (-.45,-.4) rectangle (3.15,.4);
\node[lab,anchor=west] at (3.5,0) {$|\eta(X)|=|0|=0$};
% row 2: {1,2},{3,4}
\draw[grp] (-.45,-1.5) rectangle (1.35,-.7);
\draw[grp] (1.35+.0,-1.5) rectangle (3.15,-.7);
\node[lab,anchor=west] at (3.5,-1.1) {$|{+1}|+|{-1}|=2$};
% row 3: {1,3},{2,4} 用弧线连接同组的点
\draw[FigureBlue,line width=.8pt] (a3.north) to[out=40,in=140] (c3.north);
\draw[FigureRed,line width=.8pt,dashed] (b3.south) to[out=-40,in=-140] (d3.south);
\node[lab,anchor=west] at (3.5,-2.55) {$|{+3}|+|{-3}|=6$};
\node[note,anchor=west] at (3.5,-3.05) {按正负区域分割取得上确界};
\end{tikzpicture}
